% =====================================================================
%  Conferenze, Zagarolo (On The Road), giorno 2 — Intervista 4
%  Dal tripla A all'indie: due veterani fondano Day 4 Night
%  Fonte: trascrizione "Track 18". Intervista in italiano a Christian
%  Cantamessa e Cristina, cofondatori dello studio Day 4 Night, con
%  domande dal pubblico. Per accordo iniziale non si parla del gioco in
%  produzione. Trascrizione molto rovinata nei nomi propri.
%  Etichette con prefisso "c3-4:". Nessun pacchetto aggiuntivo.
% =====================================================================

\chapter[Dal tripla A all'indie: Day 4 Night]{Dal tripla A all'indie: due veterani
fondano uno studio}
\label{cap:c3-4-day-4-night}

L'intervista della domenica pomeriggio ha due ospiti con, secondo il moderatore,
``un centinaio d'anni di esperienza in due'': uno sceneggiatore e game director che
ha lavorato a \emph{Red Dead Redemption}, e una producer cresciuta in Ubisoft
Milano. Insieme ad altri soci hanno fondato da poco uno studio indipendente, che
sta producendo il suo primo gioco. La regola stabilita in apertura:
\textbf{nessuna domanda sul gioco in uscita}, perché il team è in piena
produzione; su tutto il resto si può chiedere liberamente. Il capitolo è
organizzato così:
\begin{enumerate}
  \item gli ospiti e lo studio (sezione~\ref{sec:c3-4-ospiti});
  \item che cosa si portano dietro dal tripla A, e che cosa devono abbandonare
        (sezione~\ref{sec:c3-4-esperienza});
  \item aspettative verso il team e i molti cappelli dei fondatori
        (sezione~\ref{sec:c3-4-cappelli});
  \item il nome dello studio (sezione~\ref{sec:c3-4-nome});
  \item come nascono le storie e come si gestisce la dialettica creativa
        (sezione~\ref{sec:c3-4-storie});
  \item le domande del pubblico (sezione~\ref{sec:c3-4-domande}).
\end{enumerate}

% ---------------------------------------------------------------------
\section{Gli ospiti e lo studio}
\label{sec:c3-4-ospiti}

\begin{description}
  \item[Christian Cantamessa]\footnote{Trascritto ``Christian, Catalessa'' e
        ``signor cantanessa'', poi correttamente ``Cantamessa''.} è
        \textbf{co-game director} del nuovo studio, ne scrive la sceneggiatura ed
        è \textbf{CEO} della società negli Stati Uniti. Fa videogiochi da quando
        aveva quindici anni, nel \textbf{1991}, su \textbf{Amiga
        500}\footnote{Trascritto ``Amica 500''.}. Ha lavorato in Ubisoft a
        Milano, poi in Scozia e negli Stati Uniti per \textbf{Rockstar}, dove è
        stato tra gli autori di \textbf{\emph{Red Dead
        Redemption}}\footnote{Trascritto ``Re Re Re Revolution'', ``Re Dentro'' e
        ``Reddit''.}; più di recente era in \textbf{Microsoft}, sul progetto
        \emph{Perfect Dark}\footnote{Trascritto ``Perfidar'' e ``Super Dark''.}.
  \item[Cristina]\footnote{Il cognome non è pronunciato in modo riconoscibile.}
        è cofondatrice e amministra lo studio italiano. È entrata nei videogiochi
        nel \textbf{2008}, provenendo dalla \textbf{scrittura per cinema e
        fumetti}, ed è cresciuta in Ubisoft Milano passando da ruoli creativi
        (scrittrice, designer) al QA e poi alla \textbf{produzione}, fino a
        \emph{Mario + Rabbids Kingdom Battle}\footnote{Trascritto ``Mind Grabbit
        Skin Dumb Battle''.}.
\end{description}

Lo studio si chiama \textbf{Day 4 Night}\footnote{Trascritto ``Day Fortnite'',
``dei Fortnite'', ``The Fornight''. La ricostruzione è confermata dal racconto
sull'origine del nome (sezione~\ref{sec:c3-4-nome}).} e ha \textbf{cinque
fondatori}:
\begin{itemize}
  \item sul lato contenuti, i due direttori creativi \textbf{Christian Cantamessa}
        e \textbf{Davide Soliani}\footnote{Trascritto ``Davide Sogliani'' e
        ``Dario Sognani''; Davide Soliani è stato il creative director di
        \emph{Mario + Rabbids} in Ubisoft Milano.};
  \item sul lato amministrazione e direzione: Cristina, \textbf{Gianmarco
        Zanna}\footnote{Trascritto anche ``Marco Zanna'' e ``Tornato Zanna''.} e
        \textbf{Luca Breda}\footnote{Trascritto anche ``Luca Preda''.}, che è
        anche il ``\textbf{prototipatore selvaggio notturno}'' dello studio.
\end{itemize}
Lo studio ha una società negli Stati Uniti e uno studio in Italia.

% ---------------------------------------------------------------------
\section{Che cosa si portano dietro dal tripla A}
\label{sec:c3-4-esperienza}

\subsection{Christian: gli istinti}

Ogni progetto è a sé, quindi le \textbf{lezioni specifiche} sono difficili da
applicare, soprattutto passando dal tripla A al doppia A o all'indie. Ciò che resta
è l'\textbf{allenamento}: come in palestra, più ti alleni più diventi forte. Dopo
trent'anni di mestiere Christian sa di poter contare su una serie di
\textbf{istinti}, come quando guidando ci si accorge che il cervello fa operazioni
senza pensarci attivamente. Nel game design e nella scrittura spende meno tempo a
capire che cosa funziona e che cosa no: lo sente \textbf{istintivamente}. Poi ci
sono gli strumenti del mestiere, che col tempo diventano automatici.

\subsection{Cristina: gli automatismi da tenere e quelli da abbandonare}

Alcuni automatismi diventano naturali: a volte non serve nemmeno un foglio Excel per
capire rischi, soluzioni e dove si andrà a parare nella pianificazione. Altri,
invece, passando da realtà corporate che producono giochi tripla A (e ``quadrupla
A'') a un team più piccolo, \textbf{vanno abbandonati}, e ci vuole un po' per
iniziare a ragionare in modo più \textbf{flessibile e veloce}.

Per Cristina ``indie'' significa soprattutto \textbf{indipendenza creativa}, con team
più piccoli o più grandi. La cosa fondamentale, smettendo di ragionare da tripla A,
è la necessità di \textbf{iterare molto velocemente} e di \textbf{fallire molto più
velocemente}, per trovare subito soluzioni e strade alternative.

Ciò che invece non ha potuto portarsi dietro è l'esperienza di \textbf{essere
proprietari}: lo studio è dei cinque fondatori, quindi la responsabilità non è solo
del processo produttivo. Bisogna gestire la società, farla prosperare e tradurre la
visione creativa dei ``due talenti straordinari'' in un team \textbf{super
talentuoso ma sostenibile nel lungo termine}.

% ---------------------------------------------------------------------
\section{Il team e i molti cappelli}
\label{sec:c3-4-cappelli}

\subsection{Le aspettative verso il team}

Il moderatore chiede quali aspettative abbiano verso il team, venendo da team
enormi di altissimo livello. Cristina: il \textbf{core team} è fatto di persone con
lunga esperienza, quindi la difficoltà è stata il contrario, cioè
\textbf{disabituarle} al modo di pensare corporate e renderle più flessibili. Poi,
durante la produzione, c'è la fase di \textbf{ramp-up}, in cui si aggiungono talenti
da esperienze diverse; e chi viene dall'indie porta una \textbf{freschezza e
un'energia} vitali.

Christian racconta che lui e Davide volevano assolutamente uscire da un contesto
``iper corporate'': quando hanno fondato lo studio lui era in Microsoft (``più
corporate di così non so se esista''). Si sono licenziati ``con un certo livello di
felicità'', perché il lavoro era diventato operativamente insostenibile: la mattina
non si aveva voglia di lavorare. Ora non vede l'ora di alzarsi per lavorare con il
team italiano. C'è il desiderio di tornare a \textbf{team piccoli}, agili, più
creativi, capaci di prendersi rischi che un team di 400 o 4000 persone difficilmente
può prendere.

Detto questo, sono \textbf{molto esigenti} con i collaboratori: proprio perché sono
talenti, l'obiettivo è che \textbf{giochino al massimo del loro talento}. ``Se hai
dei campioni in squadra, non puoi tenerli in panchina''.

\subsection{I cappelli}

Christian è co-game director, sceneggiatore e CEO della società americana: deve
``indossare magliette diverse'', come aveva detto un relatore il giorno
prima\footnote{Il riferimento è al consiglio di indossare letteralmente cappelli o
magliette diverse, dato nel talk del capitolo~\ref{cap:c3-3-people-over-processes}.}.
È inevitabile, ma anche positivo: ora il destino è nelle loro mani, con un potere
decisionale che prima non avevano. Occuparsi di cose nuove, come i
\textbf{copyright} dei loghi e dei personaggi, gli fa capire esigenze di business
che poi applica al lavoro creativo.

Cristina preferisce parlare dei cappelli che ha dovuto \textbf{togliere}. I
fondatori sapevano sviluppare un gioco, dalla concezione al post-lancio; hanno
dovuto imparare ad \textbf{amministrare una società}. Con Gianmarco Zanna e Luca
Breda amministra lo studio italiano, e ha dovuto togliere il cappello da producer e
da sviluppatrice ``con le mani in pasta'' ogni giorno: contratti, consulenti,
legali, commercialisti, consulenti del lavoro, molto back office. Da persona che si
riteneva esperta si è ritrovata ``\textbf{principiante assoluta}'', ad imparare a
essere amministratrice e imprenditrice. È stata la cosa più eccitante degli ultimi
anni, e spesso si è chiesta perché non l'avesse fatto prima. Ma ci voleva un
``\textbf{allineamento di pianeti}'': le persone giuste, senza le quali non
l'avrebbe mai fatto.

% ---------------------------------------------------------------------
\section{Il nome dello studio}
\label{sec:c3-4-nome}

Cristina racconta la nascita del nome. Volevano un nome italiano ma internazionale;
il latino sembrava pretenzioso; i riferimenti all'Italia, anche culinaria, troppo
``italiani''; e i domini dei nomi che piacevano erano tutti presi. Con Christian
dall'altra parte dell'oceano si ritrovavano a scriversi su WhatsApp alle quattro di
notte. Che cosa li accomunava, oltre ai videogiochi? L'amore per il
\textbf{cinema}, e il lavorare di notte e di giorno a fusi orari opposti, spesso
nelle ore ``liminali''.

Da qui l'\textbf{effetto notte}, in inglese \emph{day for night}: la tecnica
cinematografica con cui una scena girata di giorno, in pieno sole, viene fatta
passare per notte. Sostituendo ``for'' con ``4'' è nato \textbf{Day 4 Night}.

Christian aggiunge che il nome gli piace perché è un \textbf{testamento al potere
dell'immaginazione}: l'effetto notte crea la notte dal giorno, situazioni e mondi
che non esistono. Adatto a uno studio che vuole fare giochi in cui storie,
immaginazione e creatività sono centrali.

% ---------------------------------------------------------------------
\section{Come nascono le storie}
\label{sec:c3-4-storie}

\subsection{Situazione e personaggi}

Qual è il ``segreto'' delle storie di Christian? Non c'è, risponde: è come chiedere
da dove vengono le idee, e ogni filosofo darebbe una risposta diversa. In pratica,
di solito parte da una \textbf{situazione}: per esempio ``la fine del vecchio West'',
oppure ``che cosa succederebbe se un personaggio di videogiochi uscisse dal suo
mondo e vivesse avventure fuori dal suo contesto?''. Subito dopo passa ai
\textbf{personaggi}, almeno al protagonista. Con questi due momenti iniziali si
trova di solito in una buona posizione.

\subsection{Catalogare le idee}

Christian scrive \textbf{tutte le idee}, perché spesso sono parziali: un tempo su un
taccuino, ora nelle note del telefono. Spesso questi frammenti
\textbf{si coagulano} in una macro-idea: due o tre cose apparentemente scollegate,
scritte a mesi o a un anno di distanza, si rivelano collegate, ed è l'inizio di
qualcosa. Ogni idea ha il suo \textbf{mezzo}: quelle per i videogiochi in un
``secchio'', quelle per il cinema in un altro.

\subsection{La storia è gameplay}

Christian è convinto che \textbf{la storia sia gameplay e il gameplay sia storia}.
Se ha un'idea puramente narrativa, sotto c'è un elemento di gameplay, anche se a
volte non riesce a vederlo lui: Davide, più forte sul gameplay e sulle meccaniche,
``lancia la palla'' con l'elemento di gioco. E viceversa: quando Davide propone
idee di puro gameplay, Christian sta già immaginando che cosa significano a livello
di storia. Pensare che gameplay e storia siano separati ``è una falsità''.

\subsection{La dialettica tra due direttori creativi}

Come si tengono insieme due forti personalità creative? Cristina: la creatività
nasce anche quando è più difficile. La \textbf{dialettica} non significa essere
sempre d'accordo: dal disaccordo nasce la scintilla dell'originalità. Lo studio fa
quello che chiamano \textbf{``gioco d'autore''}: gli autori sono Christian e Davide,
e il resto dello studio è al servizio di quella visione. Il compito non è
eliminare la dialettica, ma tenerla viva e farla \textbf{concretizzare in un
progetto}.

A volte la creatività va imbrigliata, ma dai \textbf{vincoli} (\emph{constraint})
un creativo trova soluzioni che ``a briglia sciolta'' non verrebbero fuori. Il
lavoro della gestione è tradurre una visione creativa forte in un progetto
produttivo, con un \textbf{equilibrio} (``non mi piace dire compromesso'') tra
fattibilità e autorialità.

Christian aggiunge due principi:
\begin{itemize}
  \item \textbf{non si litiga davanti al team}: davanti al team i due direttori sono
        completamente uniti;
  \item c'è sempre \textbf{rispetto} per le idee dell'altro; usa l'immagine della
        \textbf{partita di basket}: se Davide gli passa la palla ma nessuno dei due è
        in posizione per tirare, continuano a passarsela finché non arrivano a
        canestro.
\end{itemize}

% ---------------------------------------------------------------------
\section{Le domande del pubblico}
\label{sec:c3-4-domande}

\subsection{Che cosa portare dal percorso corporate?}

Un partecipante chiede quali scelte del loro percorso siano state le migliori, e
che cosa si possa imparare senza ripetere lo stesso percorso in grandi aziende.
Cristina: è una domanda difficile, perché hanno iniziato quando mercato e
professione erano molto diversi, ``tutto da inventare'', senza corsi né accademie.
Lei, entrata nel 2008, è stata ``gettata nell'oceano senza saper nuotare'' e ha
imparato tutto sul campo; c'erano meno sviluppatori, meno giochi, meno soldi, ma
anche meno competizione. Il loro percorso \textbf{non è replicabile}, perché sono
cambiati i tempi.

Racconta il suo ingresso: la casa editrice per cui pubblicava fumetti stava
chiudendo, e mandò il curriculum a un'offerta per \textbf{game designer} senza sapere
che cosa fosse. Non la chiamarono; un anno dopo la richiamarono perché serviva una
\textbf{scrittrice}. ``A volte bisogna buttarsi.''

Il suo consiglio: \textbf{perseverare}, ma \textbf{fare tanta esperienza} prima di
sentirsi pronti a un'avventura imprenditoriale. A vent'anni non sarebbe stata pronta:
servono ``spalle larghe''. Fare esperienza in altre aziende, fare game jam,
costruirsi un curriculum e delle demo, frequentare l'ambiente.

\subsection{Quando è il momento di aprire uno studio?}

Christian: non c'è una ricetta. Cita \textbf{David Fincher}, entrato nel cinema
attraverso i videoclip: il cinema è una \textbf{cittadella circondata da un muro};
il tuo lavoro è trovare una \textbf{crepa} ed entrare, ma chi è dentro la richiude
subito, e il prossimo deve trovarne un'altra. Vale per tutte le industrie
dell'intrattenimento. Quando ha cominciato lui non c'erano nemmeno libri su come
fare videogiochi.

Due cose però sono necessarie:
\begin{enumerate}
  \item \textbf{l'esperienza}: lanciare un'azienda senza esperienza è come voler
        fare il chirurgo senza studiare medicina. ``Se va bene per la medicina, non
        capisco perché non debba andare bene per i videogiochi.'' Che l'esperienza
        si faccia da soli (incubatori, game jam, progetti personali) o in una mega
        corporate come Microsoft, conta farla;
  \item \textbf{i crediti}: per trovare i soldi ci sono molti modi (``la lotteria,
        la famiglia ricca'', investitori stranieri, spesso stando all'estero,
        realtà di investimento italiane, banche), ma chiunque investa vuole sapere
        \emph{a chi} dà i soldi. La risposta migliore è una \textbf{lista di
        crediti}, magari con le copie vendute: il successo garantito non esiste, ma
        l'investitore ha un minimo di verifica.
\end{enumerate}
Il consiglio generico: \textbf{accumulare esperienza e crediti}. ``Se uno vince la
lotteria, può saltare entrambe le cose.''

\subsection{Quando un'idea è matura?}

Come capisce Christian quando un'idea catalogata è pronta a ``prendere il volo''?
Quando \textbf{non riesce a smettere di pensarci}. Negli ultimi quindici anni ha
preso l'abitudine di cercare deliberatamente di smettere di pensare alle idee e di
metterle da parte. Quando un'idea non lo abbandona, continua a tornare a galla (un
sogno, qualcosa visto in televisione), e viene voglia di lavorarci, è il momento
giusto. Altre idee, che sembravano bellissime appena scritte, dopo due o tre giorni
sono dimenticate: tornano ``in salamoia'' finché qualcosa non cambia.

\subsection{Quali competenze si imparano solo nelle grandi aziende?}

Cristina: le \textbf{soft skill} sono le stesse ovunque (lavorare in team, l'umiltà
di accettare le idee altrui). Ciò che insegna una corporate sono i \textbf{processi su
larga scala}: lavorare con centinaia di persone e avere una visione ampia, che poi si
traduce in realtà più piccole \textbf{semplificando i processi}. Soprattutto, si
sviluppa l'automatismo di \textbf{prevedere dove nascerà il problema}, perché lo si
è già visto molte volte, e su larga scala ogni problema diventa gravissimo. È come
nei videogiochi quando, a un certo punto, un gesto diventa automatico: non servono
più nemmeno i numeri (anche se vanno comunque prodotti), basta il riflesso.

\subsection{GTA VI è replicabile?}

Una domanda sul ``fenomeno \emph{GTA VI}'' (uscito nei giorni della
conferenza)\footnote{Si veda l'introduzione al primo giorno della conferenza.}: si
andrà verso giochi sempre più enormi, o verso giochi più piccoli?

Christian, che ha lavorato in Rockstar\footnote{Nella trascrizione ``Orchestra'',
con ogni probabilità ``Rockstar''.}, pensa che \emph{GTA VI} sia
\textbf{l'ultimo esempio} di un modo di fare videogiochi non sostenibile per
nessun'altra azienda e per nessun altro brand. Rockstar ha fatto una scommessa
relativamente sicura: \emph{GTA Online} ha finanziato parte dello sviluppo, con
incassi mensili che altri giochi non vedono; ma non esiste un altro brand per cui
un'azienda investirebbe ``un miliardo''. È il \textbf{canto del cigno} di
un'industria con enormi costi fissi e team giganteschi.

Il motivo: siamo alle soglie di \textbf{world model} e \textbf{neural rendering}, che
non sono IA generativa ``a prompt'' ma \textbf{strumenti}, che finiranno
nell'hardware delle nuove console e delle schede NVIDIA. Fisica, texture, ray
tracing e path tracing passeranno all'hardware, come negli anni Novanta con
l'\textbf{accelerazione 3D}: prima delle GPU i poligoni si disegnavano via software,
con algoritmi di riempimento scritti a mano\footnote{Trascritto
``polifilters''.}, e oggi nessuno li scrive più. Allo stesso modo non serviranno
4000 persone ``che fanno l'erba'': un team molto più piccolo potrà generare con un
world model un mondo di quella scala.

Cristina aggiunge che \emph{GTA VI} non è solo un gioco, è un \textbf{fenomeno di
costume}, un \emph{unicum}: ricoprire i grattacieli di tutte le città con la sua
pubblicità non è replicabile. In questi anni sono invece nati \textbf{generi nuovi} e
giochi più piccoli che riempiono nicchie che non si sapeva esistessero. In un mercato
difficile, il circolo virtuoso è non replicare formule note ma \textbf{dare al
pubblico ciò che non sa ancora di volere}. Replicare un genere di successo per
uscire due anni dopo, quando funzionerà già altro, non ha senso. Il pubblico ha
premiato contenuti \textbf{originali} e umani: quando sente che un progetto nasce
dalla \textbf{passione} (``termine ormai abusato''), lo abbraccia.

\subsection{Iterazione e direzione creativa}

Quanto e come iterano sulla direzione creativa? Christian: iterare è
\textbf{l'essenza} del mestiere. Fare un videogioco è come fare una
\textbf{scultura}: si prende il blocco di marmo e si toglie, si guarda, si toglie
ancora; come diceva Michelangelo, ha ``liberato'' il David dal marmo. Con un
prodotto interattivo non si può essere certi che qualcosa funzioni: \textbf{metterlo
per scritto è quasi una perdita di tempo}, serve solo a comunicare un'idea che va
realizzata almeno in un \textbf{prototipo}.

Per questo si affidano molto a Luca Breda, il cofondatore che fa i prototipi: con un
pitch brevissimo, vedono l'idea in azione dopo \textbf{24--48 ore} (``magari non
dorme per riuscirci''). Se funziona, passa a qualcun altro del team per essere
raffinata; oppure chiedono a chi è forte in un'area di sperimentare. Ma
\textbf{finché non è testabile, non esiste}. Lo stesso vale per la scrittura: il
processo non finisce con la prima stesura ma ``con la ventesima'', e continua
finché il progetto non ti viene tolto dalle mani per uscire.

\subsection{Il momento di mettersi in proprio}

Che cosa li ha portati a lasciare le grandi aziende, e qual è stato il catalizzatore?

\paragraph{Cristina.} Era una riflessione che faceva da tempo, ma aspettava che i
tempi fossero maturi, perché non è una scelta da fare da soli. In Ubisoft aveva già
assaporato che cosa volesse dire un progetto ``indie'': il primo \emph{Mario +
Rabbids} all'inizio era un progetto in cui \textbf{nessuno credeva}, cosa rarissima,
e le piaceva l'idea di ritrovare quell'ebbrezza di libertà creativa, non sua ma del
team, visto che è passata alla produzione. Precisa: i fondatori sono cinque, ma i
creativi sono due; lei e Gianmarco Zanna fanno ``un passo indietro'', il che non
significa non avere idee. Si sentiva ``arrivata'', non perché fosse soddisfatta, ma
perché lì non sarebbe potuta andare oltre, e aveva ancora forza ed entusiasmo per una
sfida nuova. Quando le persone giuste si sono allineate, hanno detto: facciamolo.

\paragraph{Christian.} Dopo l'uscita di \emph{Red Dead Redemption} si era già
licenziato da Rockstar e lavorava da freelance a progetti di videogiochi, cinema e
fumetti\footnote{La trascrizione cita qui un titolo non ricostruibile (``Shadow
Murder'').}. Il progetto gli aveva dato molto ma gli era anche costato molto sul
piano personale: \textbf{cinque anni e mezzo} di sviluppo. Da freelance era finito
a lavorare in Microsoft su \emph{Perfect Dark}, a tempo pieno, perdendo
l'indipendenza che aveva quando poteva lavorare contemporaneamente a un film e al
suo fumetto. Era una fase di sofferenza personale; la sua reazione istintiva sarebbe
stata andarsene a fare altro, magari non fare nulla per un anno.

In quel momento Davide lo chiamò per gli auguri di Capodanno. Alla domanda ``come
stai?'' Christian gli raccontò la situazione, e Davide rispose qualcosa di simile:
aveva finito un grosso progetto in Ubisoft e cercava il suo futuro, lì o altrove
(``speculazione selvaggia, no comment''). I due si conoscevano da quando avevano
lavorato insieme in Ubisoft a Milano. ``Perché non facciamo qualcosa insieme?''
Hanno cominciato a scambiarsi idee su Zoom, finché un'idea è diventata ``sarebbe
bello fare questo gioco''. Il passo successivo: ``chi ci impedisce di farlo? Stiamo
diventando vecchi: adesso o mai più''. Lo hanno fatto, ``probabilmente nel momento
peggiore dell'industria'': ``una volta premuto il pulsante di \emph{eject}, speri
che si apra il paracadute''.

\subsection{Una cultura in cui tutti contribuiscono}

Un professionista del QA in uno studio inglese osserva che, come gli ha raccontato
Davide, nello studio tutti giocano regolarmente al gioco e si osservano i nuovi
arrivati per raccogliere la loro prospettiva. Il QA, per lui, non è solo trovare i
bug ma verificare che il gioco restituisca l'esperienza prevista. Come si costruisce
una cultura in cui tutti i dipartimenti contribuiscono alla visione creativa?

Cristina (che nel passaggio da creativa a producer è passata anche dal QA) risponde
che è questione di \textbf{equilibrio e comunicazione}. Da un lato il team,
appassionato e talentuoso, deve sentirsi \textbf{coinvolto} e abbracciare la visione
dando il proprio contributo; dall'altro la visione deve restare \textbf{unica}, un
``centro di gravità permanente'', senza il quale il gioco \textbf{si sfilaccia} e
perde identità. In un team relativamente piccolo ognuno può avere un
\textbf{rapporto diretto} con la direzione creativa, senza strati di riporti: nei
grandi progetti a cui hanno lavorato, un artista poteva non avere mai la possibilità
di parlare con il direttore creativo, magari in un altro studio. L'organizzazione è
il più possibile \textbf{orizzontale}, con i livelli di management necessari a
pianificare, e si cerca di far comunicare ogni dipartimento direttamente con la
visione creativa.

Christian concorda: la comunicazione è la cosa più importante. Cita la
\textbf{metafora dell'orchestra} di \textbf{Steve Jobs} (``per quanto possa essere
antipatico''): in un'orchestra come la Los Angeles Philharmonic ogni musicista è un
virtuoso del proprio strumento, ma seduti tutti insieme \textbf{non riescono a
sentire la musica}, perché suona tutta l'orchestra. Serve qualcuno davanti, con la
partitura, che aiuti questi talenti a sentirla: volumi, attacchi, rallentamenti.
``Io non suono uno strumento, suono l'orchestra.'' Non è sminuire il talento
individuale: senza musicisti di talento non c'è musica, e lui e Davide da soli non
potrebbero farla.

\subsection{L'ispirazione dietro \emph{Red Dead Redemption}}

L'ultima domanda viene da un fan di \emph{Red Dead Redemption}: da dove è venuta
l'ispirazione, e che consigli dare a chi comincia?

\paragraph{Un western per incarico.} Christian chiarisce che in una grande azienda
c'è una componente corporate: era a \textbf{Edimburgo} quando gli dissero di
trasferirsi a \textbf{San Diego} per seguire un progetto che stava iniziando, un
western. Non si era svegliato una mattina con la voglia di fare un western. A San
Diego, parlando con il team, ha riflettuto sui possibili tipi di western e su quale
\textbf{quel team} sarebbe stato in grado di fare. Senza sminuire il lavoro di
scrittura suo e degli altri due autori, \textbf{Dan Houser} e \textbf{Michael
Unsworth}\footnote{Trascritti ``Dan Hauser'' e ``Michael Lansworth''.}, all'origine
della storia c'è stata una riflessione su ciò che il team poteva e voleva fare.

\paragraph{Il problema delle armi.} Il team voleva fare uno \textbf{shooter}. Ma nel
western classico, a fine Ottocento, ci sono solo il revolver a sei colpi, il
fucile a ripetizione e i moschetti ad avancarica: come meccaniche da shooter ``non
sono il massimo''. Era un problema logistico notevole, che spiegava la ``regola
d'oro'' secondo cui non si possono fare giochi western. Un designer appassionato di
armi gli fece notare che, guardando le date di invenzione, nella \textbf{Prima
guerra mondiale} c'erano già mitragliatrici, granate, fucili a pompa. Spostando la
data dalla fine dell'Ottocento ai \textbf{primi del Novecento, intorno al 1910}, si
ottenevano tutte quelle armi e un periodo molto interessante della storia del West:
quello che divenne la tagline, ``\textbf{The Death of the West}'', la fine del West.
Christian se la scrisse su un post-it attaccato al monitor, ``e da lì è nato tutto
il resto del gioco''. L'ispirazione, quindi, può arrivare anche dalle
\textbf{meccaniche} e da decisioni molto pragmatiche.

\paragraph{Consigli a chi comincia.} Iniziare oggi è molto più difficile di quando
ha iniziato lui, ma contano sempre \textbf{talento ed esperienza}. Il suo consiglio:
\begin{itemize}
  \item \textbf{fare cose}, fare pratica: oggi gli strumenti sono accessibili a
        tutti (motori gratuiti come Unreal\footnote{Trascritto ``Reera''.} e
        Unity, Blender per la modellazione, e anche le IA, che per imparare
        ``non vanno demonizzate'');
  \item \textbf{fare networking}: conoscere persone con gli stessi interessi,
        lavorare insieme, parlare con la community. Non si sa dove porteranno quelle
        connessioni: all'università non conta tanto ciò che insegnano quanto la
        possibilità di fare pratica e i \textbf{contatti con gli altri studenti}.
        Basta che uno di loro abbia successo e faccia partire un progetto, e potrebbe
        chiamarti: ``abbiamo fatto quel bel progetto insieme un paio d'anni fa,
        vieni a lavorare con me?''.
\end{itemize}
``Il network delle amicizie e la pratica sono le due cose più importanti.''

% ---------------------------------------------------------------------
\section*{Punti chiave del capitolo}
\addcontentsline{toc}{section}{Punti chiave del capitolo}

\begin{itemize}
  \item Dall'esperienza tripla A restano gli \textbf{istinti} e la capacità di
        prevedere i problemi; vanno abbandonati i processi pesanti, a favore di
        iterazioni e fallimenti rapidi.
  \item Fondare uno studio significa togliersi alcuni cappelli (producer,
        sviluppatore) e indossarne di nuovi (amministrazione, legale, business).
  \item ``Gioco d'autore'': una visione creativa unica, tenuta viva dalla
        dialettica tra due direttori che non litigano mai davanti al team.
  \item Le storie nascono da una \textbf{situazione} e dai personaggi; le idee si
        catalogano e maturano quando non smettono di tornare; storia e gameplay non
        sono separabili.
  \item Iterare è l'essenza del mestiere: un'idea non testabile non esiste; si
        prototipa in 24--48 ore.
  \item Per aprire uno studio servono \textbf{esperienza} e \textbf{crediti}; per
        entrare nell'industria, \textbf{pratica} e \textbf{networking}.
  \item \emph{GTA VI} come canto del cigno dei mega-team: world model e neural
        rendering potrebbero ridurre drasticamente le dimensioni dei team.
  \item \emph{Red Dead Redemption}: l'ambientazione nel 1910 nacque da un vincolo di
        gameplay (le armi), da cui la tagline ``The Death of the West''.
\end{itemize}
